% META: {"id": "TCOMPL-MF-EX-009", "chapitre": "TCOMPL-MODELES-FONCTION", "type_objet": "exercice", "capacites": ["TCOMPL-MF-C5"], "capacites_codes": ["C5"], "methodes": ["M5"], "parcours": 2, "competences": ["calculer", "raisonner"], "duree_min": 15, "mode_creation": "ex_nihilo", "sources_inspiration": [], "fichier_tex": "chapitres/TCOMPL-MODELES-FONCTION/exercices/TCOMPL-MF-EX-009.tex", "corrige_tex": "chapitres/TCOMPL-MODELES-FONCTION/corriges/TCOMPL-MF-CO-009.tex", "status": "generated"}
% BEGIN-VERIFY
% from sympy import symbols, diff, exp, simplify, solve
% x = symbols('x', real=True)
% f = x * exp(-x)
% assert simplify(diff(f, x) - (1 - x) * exp(-x)) == 0
% assert simplify(diff(f, x, 2) - (x - 2) * exp(-x)) == 0
% assert solve(diff(f, x, 2), x) == [2]
% assert diff(f, x, 2).subs(x, 0) < 0
% assert diff(f, x, 2).subs(x, 3) > 0
% assert simplify(f.subs(x, 2) - 2 * exp(-2)) == 0
% END-VERIFY
\begin{exercice}{TCOMPL-MF-EX-009}{2}{15}
Soit $f(x)=x\,\mathrm{e}^{-x}$ sur $\mathbb{R}$.
\begin{enumerate}
  \item Calculer $f'$ et $f''$.
  \item Étudier la convexité de $f$ et donner les points d'inflexion.
  \item En déduire l'allure de la courbe au voisinage de son point
        d'inflexion.
\end{enumerate}
\end{exercice}
